\scanpage{156}
\mspage{V.28}
% archive sheet number: 155

b) Les résultats précédents suggèrent la conjecture suivante \inserted{(au moins pour $A,B$ \uline{excellents})}

Soit $f:X\inserted{=\operatorname{Spec}B}\unsure{-|\longrightarrow}S\inserted{=\operatorname{Spec}A}$ un morphisme \uline{plat} de schémas noethériens, \inserted{stricts locaux, $s$ le point fermé de $S$. $X_0=f^{-1}(s)$.} \deleted{\unsure{et à les|et à la} fibres géométriques \unsure{satisfaisant|satisfassant} à la condition $(R_d)$,} \deleted{Soit \deleted{$X_0$} $k$ le corps \unsure{rés. du pt} \illegible{1},} Supposons que \deleted{$X_0$ soit géométriquement $(R_d)$, et que} pour tt extension finie $k'$ de $k=k(s)$, l'ens. \deleted{$m=\operatorname{prof\,géom}$}\inserted{singulier de} $X_0\otimes_k k'$ \deleted{satisfasse la condition $R_{\unsure{1|d}}$, et supposons} \inserted{soit de dim $\leq d$} et $m=\operatorname{prof\,géom}B$. Alors $f$ est loc. $(m-d\unsure{-|+}1)$-acyclique, \deleted{et si $m-d\unsure{-|+}1\geq\deleted{\unsure{3}}1$} pour $\unsure{\mathbb{L}|L}=\mathbb{P}-\{p\}$, ($p$. car résiduelle de $A$), et si $m-d-1\geq1$ i.e. $m-d\geq2$, $f$ est loc. $1$-asphérique pour $\mathbb{L}$. \deleted{Cela} Il est plausible que cela soit démontrable comme conséquence de la conjecture déjà citée SGA 1 $\overline{\underline{\mathrm{XIV}}}$ 6.3, et de la résolution des singularités pour les schémas \unsure{finis} sur $\unsure{S|\eta}$,
